\documentclass{article}

% packages
\usepackage{amsmath, amsthm, thmtools, amsfonts, amssymb, luacode, catchfile, tikzducks, hyperref, ifthen}
\ifcsname c@kobocompile\endcsname
	\usepackage[a5paper, total={1072pt, 1448pt}, margin=10pt, includeheadfoot]{geometry} % set page margins
\else
	\usepackage[a4paper, margin=50pt, includeheadfoot]{geometry}
\fi
\usepackage[shortlabels]{enumitem}
\usepackage[skip=3pt, indent=0pt]{parskip}

% language
\usepackage[bidi=basic, layout=tabular, provide=*]{babel}
\ifcsname c@english\endcsname
	\babelprovide[main, import]{english}
\else
	\babelprovide[main, import]{hebrew}
	\babelprovide{rl}
\fi
%\babelfont{rm}{Libertinus Serif}
\babelfont{rm}[Renderer=Harfbuzz]{Libertinus Serif}
\babelfont{sf}{Libertinus Sans}
\babelfont{tt}{Libertinus Mono}

% style
\AddToHook{cmd/section/before}{\clearpage}	% Add line break before section
\linespread{1.3}
\setcounter{secnumdepth}{0}		% Remove default number tags from sections, this won't do well with theorems
\AtBeginDocument{\setlength{\belowdisplayskip}{3pt}}
\AtBeginDocument{\setlength{\abovedisplayskip}{3pt}}
\graphicspath{ {../images/} }

% operators
\DeclareMathOperator\cis{cis}
\DeclareMathOperator\Sp{Sp}
\DeclareMathOperator\tr{tr}
\DeclareMathOperator\im{Im}
\DeclareMathOperator\re{Re}
\DeclareMathOperator\diag{diag}
\DeclareMathOperator*\lowlim{\underline{lim}}
\DeclareMathOperator*\uplim{\overline{lim}}
\DeclareMathOperator\rng{rng}
\DeclareMathOperator\Sym{Sym}
\DeclareMathOperator\Arg{Arg}
\DeclareMathOperator\Log{Log}
\DeclareMathOperator\dom{dom}
\DeclareMathOperator\supp{Supp}
\DeclareMathOperator\var{Var}
\DeclareMathOperator\cov{Cov}

% commands
%\renewcommand\qedsymbol{\textbf{מש''ל}}
%\renewcommand\qedsymbol{\fbox{\emoji{lizard}}}
\newcommand{\Aa}[0]{\mathcal{A}}
\newcommand{\Bb}[0]{\mathcal{B}}
\newcommand{\CC}[0]{\mathbb{C}}
\newcommand{\Cc}[0]{\mathcal{C}}
\newcommand{\EE}[0]{\mathbb{E}}
\newcommand{\FF}[0]{\mathbb{F}}
\newcommand{\Ff}[0]{\mathcal{F}}
\newcommand{\Ii}[0]{\mathcal{I}}
\newcommand{\Gg}[0]{\mathcal{G}}
\newcommand{\Ll}[0]{\mathcal{L}}
\newcommand{\Mm}[0]{\mathcal{M}}
\newcommand{\NN}[0]{\mathbb{N}}
\newcommand{\Nn}[0]{\mathcal{N}}
\newcommand{\PP}[0]{\mathbb{P}}
\newcommand{\Pp}[0]{\mathcal{P}}
\newcommand{\QQ}[0]{\mathbb{Q}}
\newcommand{\RR}[0]{\mathbb{R}}
\newcommand{\Rr}[0]{\mathcal{R}}
\newcommand{\Ss}[0]{\mathcal{S}}
\newcommand{\TT}[0]{\mathbb{T}}
\newcommand{\Uu}[0]{\mathcal{U}}
\newcommand{\Vv}[0]{\mathcal{V}}
\newcommand{\Ww}[0]{\mathcal{W}}
\newcommand{\ZZ}[0]{\mathbb{Z}}
\newcommand{\acts}[0]{\circlearrowright}
\newcommand{\explain}[2] {
	\begin{flalign*}
		 && \text{#2} && \text{#1}
	\end{flalign*}
}
\newcommand{\maketitleprint}[0]{ \begin{center}
	%\begin{tikzpicture}[scale=3]
	%	\duck[graduate=gray!20!black, tassel=red!70!black]
	%\end{tikzpicture}	
	\includegraphics[width=6cm]{cover}
\end{center}
}

% theorem commands
\newtheoremstyle{c_remark}
	{}	% Space above
	{}	% Space below
	{}% Body font
	{}	% Indent amount
	{\bfseries}	% Theorem head font
	{}	% Punctuation after theorem head
	{.5em}	% Space after theorem head
	{\thmname{#1}\thmnumber{ #2}\thmnote{ \normalfont{\text{(#3)}}}}	% head content
\newtheoremstyle{c_definition}
	{3pt}	% Space above
	{3pt}	% Space below
	{}% Body font
	{}	% Indent amount
	{\bfseries}	% Theorem head font
	{}	% Punctuation after theorem head
	{.5em}	% Space after theorem head
	{\thmname{#1}\thmnumber{ #2}\thmnote{ \normalfont{\text{(#3)}}}}	% head content
\newtheoremstyle{c_plain}
	{3pt}	% Space above
	{3pt}	% Space below
	{\itshape}% Body font
	{}	% Indent amount
	{\bfseries}	% Theorem head font
	{}	% Punctuation after theorem head
	{.5em}	% Space after theorem head
	{\thmname{#1}\thmnumber{ #2}\thmnote{ \text{(#3)}}}	% head content

\ifcsname c@english\endcsname
	\theoremstyle{plain}
	\newtheorem{theorem}{Theorem}[section]
	\newtheorem{lemma}[theorem]{Lemma}
	\newtheorem{proposition}[theorem]{Proposition}
	\newtheorem*{proposition*}{Proposition}
	%\newtheorem{corollary}[theorem]{אין חלופה עברית}

	\theoremstyle{definition}
	\newtheorem{definition}[theorem]{Definition}
	\newtheorem*{definition*}{Definition}
	\newtheorem{example}{Example}[section]
	\newtheorem{exercise}{Exercise}[section]

	\theoremstyle{remark}
	\newtheorem*{remark}{Remark}
	\newtheorem*{solution}{Solution}
	\newtheorem{conclusion}[theorem]{Conclusion}
	\newtheorem{notation}[theorem]{Notation}
\else
	\theoremstyle{c_plain}
	\newtheorem{theorem}{משפט}[section]
	\newtheorem{lemma}[theorem]{למה}
	\newtheorem{proposition}[theorem]{טענה}
	\newtheorem*{proposition*}{טענה}
	%\newtheorem{corollary}[theorem]{אין חלופה עברית}

	\theoremstyle{c_definition}
	\newtheorem{definition}[theorem]{הגדרה}
	\newtheorem*{definition*}{הגדרה}
	\newtheorem{example}{דוגמה}[section]
	\newtheorem{exercise}{תרגיל}[section]

	\theoremstyle{c_remark}
	\newtheorem*{remark}{הערה}
	\newtheorem*{solution}{פתרון}
	\newtheorem{conclusion}[theorem]{מסקנה}
	\newtheorem{notation}[theorem]{סימון}
\fi

% Questions related commands
\newcounter{question}
\setcounter{question}{1}
\newcounter{sub_question}
\setcounter{sub_question}{1}

\ifcsname c@english\endcsname
	\newcommand{\question}[1][0]{
		\ifthenelse{#1 = 0}{}{\setcounter{question}{#1}}
		\section{Question \arabic{question}}
		\addtocounter{question}{1}
		\setcounter{sub_question}{1}
	}

	\newcommand{\subquestion}[1][0]{
		\ifthenelse{#1 = 0}{}{\setcounter{sub_question}{#1}}
		\subsection{Part \alph{sub_question}}
		\addtocounter{sub_question}{1}
	}
\else
	\newcommand{\question}[1][0]{
		\ifthenelse{#1 = 0}{}{\setcounter{question}{#1}}
		\section{שאלה \arabic{question}}
		\addtocounter{question}{1}
		\setcounter{sub_question}{1}
	}

	\newcommand{\subquestion}[1][0]{
		\ifthenelse{#1 = 0}{}{\setcounter{sub_question}{#1}}
		\subsection{סעיף \localecounter{letters.gershayim}{sub_question}}
		\addtocounter{sub_question}{1}
	}
\fi

% import lua and start of document
\directlua{common = require ('../common')}

\GetEnv{AUTHOR}

% headers
\author{\AUTHOR}
\date\today

\title{פתרון מטלה 05 --- אנליזה פונקציונלית, 80417}

\begin{document}
\maketitle
\maketitleprint[teal]

\question{}
יהי $T : C[0, 1] \to \RR$ פונקציונל מתמטי (כלומר העתקה לינארית) רציף ונניח שלכל $k \in \NN \cup \{ 0 \}$ מתקיים,
\[
	T(x^k)
	= \frac{1}{k + 1}
\]
נראה שלכל $f \in C[0, 1]$ מתקיים,
\[
	T(f)
	= \int_{0}^{1} f(x)\ dx
\]
\begin{proof}
	ממשפט הקירוב של ויירשטראס אנחנו יודעים שקיימים ${\{ p_n \}}_{n = 1}^\infty \subseteq C[0, 1]$ סדרת פולינומים המתכנסים במידה שווה ל־$f$. \\
	אנו גם יודעים שאופרטורים לינאריים משמרים התכנסות במידה שווה, ולכן $T(p_n) \rightrightarrows T(f)$.
	אילו נניח ש־$p_n(x) = \sum_{i = 0}^N \alpha_n^i x^i$ אז נובע,
	\[
		T(p_n)
		= \sum_{i = 0}^N \alpha_n^i \cdot \frac{1}{k + 1}
		= \int_{0}^{1} \alpha_n^i \cdot x^i
	\]
	ולכן נובע ש־$T(p_n) = \int_{0}^{1} p_n(x)\ dx$ לכל $n \in \NN$.
	אבל אינטגרל משמר התכנסות במידה שווה ונובע,
	\[
		T(f)
		= \lim_{n \to \infty} T(p_n)
		= \lim_{n \to \infty} \int_{0}^{1} p_n(x)\ dx
		= \int_{0}^{1} f(x)\ dx
	\]
	ונסיק שהטענה אכן מתקיימת.
\end{proof}

\question{}
עבור $k \in \NN$,
נראה שאוסף הפולינומים צפוף במרחב $C^k[0, 1]$ בנורמה על $C^k$ המוגדרת על־ידי,
\[
	\lVert f \rVert_k
	= \sum_{j = 0}^k \lVert f^{(j)} - p_n^{(j)} \rVert_\infty
\]
\begin{proof}
	נוכיח את הטענה באינדוקציה על $k$.
	עבור $k = 0$ הטענה מתלכדת עם משפט הקירוב של ויירשטראס ולכן סיימנו.

	נניח שהטענה נכונה עבור $k - 1$ ונבדוק את הטענה עבור $k$. \\
	נגדיר ${\{ p_n \}}_{n = 1}^\infty \subseteq C^k[0, 1]$ סדרת פולינומים כך ש־$p_n \rightrightarrows f^{(k)}$.
	נגדיר את הסדרה ${\{ p_n^1 \}}_{n = 1}^\infty \subseteq C^k[0, 1]$ על־ידי,
	\[
		p_n^1(x) = f^{(k - 1)}(0) + \int_{0}^{x} p(x)\ dx
	\]
	אנו יודעים כי לכל פולינום אינסוף קדומות וכי אף הן פולינומים, וכן ש־$\frac{d}{dx} p_n^1 = p_n$ לכל $n \in \NN$.
	מטענה מתרגול 4 נובע,
	\[
		p_n^1
		\rightrightarrows f^{(k - 1)}(0) + \int_0^x f^{(k)}(x)\ dx
		= f^{(k - 1)}(x) + f^{k - 1}(0) - f^{k - 1}(0)
		= f^{(k - 1)}(x)
	\]
	אבל מהנחת האינדוקציה נובע ש־$p_n^i \rightrightarrows f^{(i)}$ לכל $0 \le i < k$ וסיימנו.
\end{proof}

\question{}
במרחב $C[0, 1]$ עם נורמת $\lVert \cdot \rVert_\infty$ נגדיר $W = \Sp\{ x, x^2, \ldots \}$ ו־$D = \{ f \in C[0, 1] \mid f(0) = 0 \}$. \\
נראה ש־$\overline{W} = D$.
\begin{proof}
	תהי סדרת פולינומים מתכנסת במידה שווה ${\{ p_n \}}_{n = 1}^\infty \subseteq W$, ונסמן $p_n \rightrightarrows f$, אז $\lim_{n \to \infty} p_n(0) = f(0)$.
	אבל אנו יודעים כי $x \mid p_n$ לכל $n \in \NN$ מהגדרת $W$ ולכן $p_n(0) = 0$ בלבד, ונסיק $f(0) = 0$, כלומר $f \in D$, ונקבל $\overline{W} \subseteq D$.

	ממשפט הקירוב של ויירשטראס אנו יודעים כבר כי קבוצת הפולינומים צפופה ב־$C[0, 1]$.
	תהי $f \in D$ ותהי ${\{ p_n \}}_{n = 1}^\infty \subseteq C[0, 1]$ סדרת פולינומים כך ש־$p_n \rightrightarrows f$ ממשפט הקירוב.
	נגדיר גם את סדרת הפונקציות ${\{ g_n \}}_{n = 1}^\infty \subseteq C[0, 1]$ על־ידי $g_n(x) = p_n(0)$.
	אנו יודעים ש־$p_n(0) \to f(0) = 0$ ולכן הסדרה $\{ g_n \}$ מתכנסת במידה שווה, ובהתאם גם $\{ p_n - g_n \}$ היא סדרת פולינומים מתכנסת במידה שווה.
	נסיק ש־$p_n - g_n \rightrightarrows f$, ונשאר להראות ש־$\{ p_n - g_n \} \subseteq D$.
	יהי $n \in \NN$ ונניח ש־$p_n(x) = \sum_{n = 0}^N \alpha_i^n x^n$, אז $g_n(x) = \alpha_0^n$ ובהתאם $(p_n - g_n)(x) = \sum_{i = 1}^N \alpha_i^n x^n$ ולכן $p_n \in W$.
	כלומר $f \in \overline{W}$.

	מצאנו אם כך ש־$\overline{W} = D$.
\end{proof}

\question{}
יהי $(K, \rho)$ מרחב מטרי קומפקטי.
נראה ש־$\operatorname{Lip}(K)$ צפופה ב־$C(K)$.
\begin{proof}
	נזכיר כי,
	\[
		\operatorname{Lip}(K)
		= \{ f \in C(K, \RR) \mid \exists L > 0, \forall x, y \in X, |f(x) - f(y)| \le L \cdot \rho(x, y) \}
	\]
	וראינו בתרגול כי זוהי אלגברה של פונקציות רציפות, אנו נראה כי היא אף מפרידה נקודות ולא מתאפסת ב־$K$.

	נניח ש־$x, y \in K$ וכן ש־$x \ne y$, ונרצה להראות שקיימת $f \in \operatorname{Lip}(K)$ כך ש־$f(x) \ne f(y)$.
	נגדיר $f(t) = \rho(t, x)$ לכל $t \in K$.
	באינפי 3 ראינו כי פונקציה זו רציפה, וכמובן מהגדרתה היא $1$־ליפשיצית, כלומר $f \in \operatorname{Lip}(K)$.
	אבל מהגדרת המטריקה $f(x) = \rho(x, x) = 0$ אבל $x \ne y \iff \rho(x, y) \ne 0 \iff f(y) \ne 0$, ולכן $f$ מפרידה את $x, y$.

	נראה כי $\operatorname{Lip}(K)$ לא מתאפסת ב־$K$.
	נגדיר את $f(x) = 1$, זוהי פונקציה קבועה ולכן $0$־ליפשיצית, אבל לכל $x \in K$ גם $f(x) \ne 0$, ולכן מעידה ש־$\operatorname{Lip}(K)$ לא מתאפסת ב־$K$.

	נסיק ממשפט סטון־ויירשטראס כי $\operatorname{Lip}(K)$ צפופה ב־$C(K)$.
\end{proof}

\end{document}
